\documentclass[10pt,a4paper]{amsart}
\usepackage[margin=19mm]{geometry}
\usepackage{amsmath,amssymb,mathtools}
\usepackage[scr=rsfs,cal=euler]{mathalfa}
\usepackage{xfrac}
\usepackage[unicode,hidelinks,bookmarks=false]{hyperref}
\newcommand{\ie}{i.e.\ }
\newcommand{\cf}{cf.\ }
\newcommand{\resp}{respectively\ }
\newcommand{\Subsection}{\subsection}
\providecommand{\nobreakdash}{\nobreak\hbox{-}}
\setlength{\emergencystretch}{2em}
\multlinegap0pt
\usepackage{bm}
\usepackage{bbm}
\usepackage{dsfont}
\usepackage{xspace}
\usepackage{cancel}
\usepackage{mathtools}
\usepackage[shortlabels]{enumitem}
\usepackage{amsthm}
\makeatletter
\@ifundefined{theorem}{\newtheorem{theorem}{Theorem}}{}
\@ifundefined{proposition}{\newtheorem{proposition}[theorem]{Proposition}}{}
\makeatother
\newcommand{\N}{\mathbf{N}}
\newcommand{\R}{\mathbf{R}}
\newcommand{\conv}{\operatorname{conv}}
\newcommand{\ocone}[2]{\operatorname{Cone}^-(#1,#2)}
\title[A construction of Kakeya Sets in Arbitrary Dimension]{A construction of Kakeya Sets in Arbitrary Dimension}
\author{Rafael J. Fern\'andez-Delgado, Mikael de la Salle}
\date{Source: arXiv:2607.14824v3 (CC BY 4.0)}
\begin{document}
\maketitle

\noindent\emph{Source and licence.} A construction of Kakeya Sets in Arbitrary Dimension. Version arXiv:2607.14824v3 (CC BY 4.0), distributed under the Creative Commons Attribution 4.0 International licence, as recorded in the arXiv licence metadata for this exact version and shown on the abstract page https://arxiv.org/abs/2607.14824v3. Full rights notice: licenses/arxiv-2607.14824-CC-BY-4.0.txt.\par\smallskip

\noindent\textit{Editorial scope.} Counted unit: the Proposition whose source label is \texttt{prop: perron tree} in the article (source0.tex lines 506--515, source label \texttt{prop: perron tree}), delivered together with the article's own complete original proof of it (source0.tex lines 516--581, one \texttt{proof} environment of 66 lines). No mathematics is written, repaired or re-derived: statement and proof are the article's own text line for line. Required context retained uncounted: eq:assumption\_shrinking (lines 316--318); prop: base case (lines 473--483). Every definition, display and label of those lines is retained unchanged. Omitted from this package and not redistributed: 1--315, 319--472, 484--505, 582--1134. Nothing inside a retained excerpt is omitted or altered: every retained body is delivered exactly as the article's lines are written, so no omission receipt of any kind accompanies this package. Citations: the retained excerpts carry no citation command at all, so no bibliography entry of the article is transcribed and no reference is redistributed. Reference adaptation: none was needed and none was made; every reference inside the delivered unit resolves inside it, so no unresolved reference or citation is produced. Printed slips kept literal: no printed word, symbol, hypothesis or number of the counted statement or of its proof was altered. Layout and class: the article's own class is \texttt{amsart} with packages including babel, geometry, float, comment, xcolor, xurl, enumitem, graphicx, subcaption, amsmath, amssymb, amsthm, mathtools, hyperref; the wrapper is the lane's pinned \texttt{amsart} editorial wrapper at 10pt on A4, which restates the theorem-like environments the retained text needs and adds the article's own macro definitions that the retained text needs (\texttt{\textbackslash N}, \texttt{\textbackslash R}, \texttt{\textbackslash conv}, \texttt{\textbackslash ocone}), with extra wrapper packages (bm, bbm, dsfont, xspace, cancel, mathtools, enumitem). The delivered numbering is the unit's own. Assets: no retained span contains a figure environment, a graphics-inclusion command, a PDF-page inclusion, a table or a TikZ picture; no image file is copied into this package, the package declares no asset dependency and no image file is redistributed with it. Licence: version arXiv:2607.14824v3 of this article is distributed under the Creative Commons Attribution 4.0 International licence, as recorded in the arXiv licence metadata for this exact version and shown on the abstract page https://arxiv.org/abs/2607.14824v3. A full rights notice is delivered at licenses/arxiv-2607.14824-CC-BY-4.0.txt. Characters: every retained excerpt is pure ASCII, so the delivered engine, XeLaTeX with its default Latin Modern text font, reports no missing character.
\begin{equation}\label{eq:assumption_shrinking} 
\boxed{\forall i\in I, \;\; \exists \;\mathbf{x}_i \in \Sigma,\;\;\Sigma_i \subset \mathcal{H}_{\mathbf{x}_i}^c(\Sigma)}
\end{equation}

\begin{proposition}\label{prop: base case}
Let $\Sigma \subset \R^d$ be a  $(d-1)$-dimensional polytope with a partition $(\Sigma_i)_{i \in I}$ satisfying \eqref{eq:assumption_shrinking}. 

$ \forall\;\mathbf{v}\notin \operatorname{Aff}(\Sigma)$, $\forall\;\mathbf{o}\in \operatorname{Aff}(\Sigma)$ and $\forall\;t\in (0,1-c]$, if $T=\operatorname{conv}(\Sigma, \mathbf{v})$, $T_i=\operatorname{conv}(\Sigma_i, \mathbf{v})$, $\widetilde{T}=\ocone{\Sigma}{\mathbf{v}}$ and $\widetilde{T}_i=\ocone{\Sigma_i}{\mathbf{v}}$, then there exist vectors $(\mathbf{w}_i)_{i\in I}\subset \overrightarrow{\operatorname{Aff}(\Sigma)}$ such that:
\begin{enumerate}[(i)]
    \item $\displaystyle \bigg|\bigg[\bigcup_{i\in I} T_i+\mathbf{w}_i\bigg]\setminus \mathcal{H}^{1-t}_\mathbf{o}(T)\bigg|\leq \frac{t^d|T|}{(1-c)^d} $.
    \item $\displaystyle \big(\widetilde{T}_i+\mathbf{w}_i\big)_{i\in I}$ are contained in $\mathcal{H}_\mathbf{o}^{1-t}(\widetilde{T})$, and have pairwise disjoint interiors if the family $(\Sigma_i,\mathbf{x}_i)_{i \in I}$ is affinely separated.
    \item The $\Sigma_i+
    \mathbf{w}_i$ are contained in $\mathcal{H}_\mathbf{o}^{1-t}(\Sigma)$ $\forall\;i$.
\end{enumerate}
\end{proposition}

\begin{proposition}\label{prop: perron tree}

Let $\mathbf{v}\notin \operatorname{Aff}(\Sigma)$, $\mathbf{o}\in \operatorname{Aff}(\Sigma)$ and $t_1, t_2,\dots,t_n\in (0,1-c]$. Write  $P_m=\prod_{j=1}^{m}(1-t_j)$ for $m\in \N$, $T=\operatorname{conv}(\Sigma, \mathbf{v})$, $T_i=\operatorname{conv}(\Sigma_i^{(n)}, \mathbf{v})$, $\widetilde{T}=\ocone{\Sigma}{\mathbf{v}}$ and $\widetilde{T}_i=\ocone{\Sigma_i^{(n)}}{\mathbf{v}}$. Then there exists vectors $(\mathbf{w}_i)_{i\in I^n}\subset \overrightarrow{\operatorname{Aff}(\Sigma)}$ such that:
\begin{enumerate}[(i)]
    \item $\displaystyle \bigg|\bigg[\bigcup_{i\in I^n} T_i+\mathbf{w}_i\bigg]\setminus\mathcal{H}_\mathbf{o}^{P_n}(T)\bigg|\leq \frac{|T|}{(1-c)^d}\sum_{i=1}^n\bigg(\prod_{j=1}^{i-1}(1-t_j)^d\bigg)t_{i}^d$.
    \item $\displaystyle \big(\widetilde{T}_i+\mathbf{w}_i\big)_{i\in I^n}$ are contained in $\mathcal{H}_\mathbf{o}^{P_n}(\widetilde{T})$, and have pairwise disjoint interiors if the family $(\Sigma_i,\mathbf{x}_i)_{i \in I}$ is affinely separated.
\item The $\Sigma_i^{(n)}+
    \mathbf{w}_i$ are contained in $\mathcal{H}_\mathbf{o}^{P_n}(\Sigma)$ $\forall\;i$.
\end{enumerate}
\end{proposition}

\begin{proof}
   We will proceed by induction in the number of iterations $n\in \N$. The base case $n=1$ is exactly Proposition~\ref{prop: base case}. Assume that the statement holds up to $(n-1)$, and let us prove it for $n$:
   
  Apply the proposition for $n=1$, to $\Sigma$, $\mathbf{v}$, $\mathbf{o}$ and $t_n$ to obtain vectors $(\mathbf{w}_{i_n}')_{i_{n}\in I}$ that satisfy, if  $T_{i_n} = \conv(\Sigma_{i_n},\mathbf{v})$ and $\widetilde{T}_{i_n}=\ocone{\Sigma_{i_n}}{ \mathbf{v}}$, the following conclusions:
\begin{enumerate}[(a)]
    \item\label{item:induction11} $\displaystyle \bigg|\bigcup_{i_n\in I} \bigg(T_{i_n}+\mathbf{w}_{i_n}'\bigg)\setminus \mathcal{H}_\mathbf{o}^{1-t_n}(T)\bigg|\leq \frac{|T|t_n^d}{(1-c)^d}$,
    \item\label{item:induction12} $\displaystyle \big(\widetilde{T}_{i_n}+\mathbf{w}_{i_n}'\big)_{i_n\in I}$ are contained in $\mathcal{H}^{1-t_n}(\widetilde{T})$, and have pairwise disjoint interiors if the family $(\Sigma_i,\mathbf{x}_i)_{i \in I}$ is affinely separated.
    \item\label{item:induction13} $\Sigma_{i_n}+\mathbf{w}_{i_n}'$ are contained in $\mathcal{H}^{1-t_n}_\mathbf{o}(\Sigma)$ $\forall \;i_n\in I$.
\end{enumerate} 

Let us re-index the $n$-iterated partition $\big(\Sigma_i^{(n)}\big)_{i \in I^n}$ as $\big(\Sigma_{i'}^{i_n}\big)_{i'\in I^{n-1},i_n \in I}$ in such a way that for every fixed $i_{n}\in I$, $\big(\Sigma_{i'}^{i_n}\big)_{i'\in I^{n-1}}$ is an $(n-1)$-iterated partition of $\Sigma_{i_n}$. By the induction hypothesis we can apply the proposition for $(n-1)$, to the simplex $\Sigma_{i_{n}}$, with $\mathbf{v}$, $\mathbf{o}$ and $t_1$, ..., $t_{n-1}\in (0,1-c]$ to obtain vectors $\big(\mathbf{w}_{i'}^{i_n}\big)_{i'\in I^{n-1}}\subset \overrightarrow{\operatorname{Aff}(\Sigma_{i_n})}=\overrightarrow{\operatorname{Aff}(\Sigma)}$ such that:
\begin{enumerate}[(a),resume]
    \item\label{item:inductionn1} $\displaystyle \bigg|\bigcup_{i'\in I^{n-1}} \bigg(\operatorname{conv}(\Sigma_{i'}^{i_n}, \mathbf{v})+\mathbf{w}_{i'}^{i_n}\bigg)\setminus \mathcal{H}_\mathbf{o}^{P_{n-1}}(T_{i_n})\bigg|\leq \frac{|T_{i_n}|}{(1-c)^d}\sum_{k=1}^{n-1} P_{k-1}^dt_k^d$,
    \item\label{item:inductionn2} $\displaystyle \big(\ocone{\Sigma^{i_n}_{i'}}{ \mathbf{v}}+\mathbf{w}^{i_n}_{i'}\big)_{i'\in I^{n-1}}$ are contained in $\mathcal{H}^{P_{n-1}}_\mathbf{o}(\widetilde{T}_{i_n})$, and have pairwise disjoint interiors if the family $(\Sigma_i,\mathbf{x}_i)_{i \in I}$ is affinely separated.
    \item\label{item:inductionn3} The $\Sigma_{i'}^{i_n}+
    \mathbf{w}^{i_n}_{i'}$ are contained in $\mathcal{H}^{P_{n-1}}_\mathbf{o}(\Sigma_{i_n})$ $\forall \;i'\in I^{n-1}$.
\end{enumerate}

With all the above, define for $i=(i',i_n)\in I^n=I\times I^{n-1}$:
\begin{equation*}
 \mathbf{w}_i=\mathbf{w}^{i_n}_{i'}+P_{n-1}\mathbf{w}_{i_n}'\in \overrightarrow{\operatorname{Aff}(\Sigma)}\qquad.
\end{equation*}
Let us check that $(\mathbf{w}_i)_{i \in I^n}$ satisfy the conclusion of the proposition. 

In order to prove property (i), using the following inclusion, valid for any sets,
\[ \bigcup_{i_n} (A_{i_n} \setminus C) \subset \Big(\bigcup_{i_n} (A_{i_n} \setminus B_{i_n})\Big) \cup \Big(\bigcup_{i_n} (B_{i_n} \setminus C)\Big),\]
to $A_{i_n} = \bigcup_{i' \in I^{n-1}} (T_{i',i_n}+\mathbf{w}_{i'}^{i_n}+P_{n-1}\mathbf{w}_{i_n}')$, $B_{i_n} = \mathcal{H}_{\mathbf{o}}^{P_{n-1}}(T_{i_n} + \mathbf{w}_{i_n}')$ and $C= \mathcal{H}_{\mathbf{o}}^{P_{n}}(T)$, we obtain
\begin{multline*}
    \Big| \bigcup_i (T_i + \mathbf{w}_i)\setminus \mathcal{H}_{\mathbf{o}}^{P_n}(T) \Big| \leq 
\sum_{i_n} \Big|\bigcup_{i'} (T_{i',i_n} + \mathbf{w}_{i'}^{i_n} + P_{n-1} \mathbf{w}_{i_n}')\setminus \mathcal{H}_{\mathbf{o}}^{P_{n-1}}(T_{i_n} + \mathbf{w}_{i_n}')\Big|\\
+ \Big| \bigcup_{i_n} \mathcal{H}_{\mathbf{o}}^{P_{n-1}}\big((T_{i_n} + \mathbf{w}_{i_n}') \setminus \mathcal{H}_{\mathbf{o}}^{1-t_n}(T)\big)\Big|.
\end{multline*} 
We have $\mathcal{H}_{\mathbf{o}}^{P_{n-1}}(T_{i_n} + \mathbf{w}_{i_n}') = \mathcal{H}_{\mathbf{o}}^{P_{n-1}}(T_{i_n}) + P_{n-1}\mathbf{w}_{i_n}'$, so by the translation-invariance of the Lebesgue measure and \ref{item:inductionn1} the first term is
\[\sum_{i_n} \Big|\bigcup_{i'} (T_{i',i_n} + \mathbf{w}_{i'}^{i_n})\setminus \mathcal{H}_{\mathbf{o}}^{P_{n-1}}(T_{i_n})\Big| \leq \sum_{i_n\in I} \frac{ |T_{i_n}|}{(1-c)^d}\sum_{k=1}^{n-1} P_{k-1}^dt_k^d = \frac{ |T|}{(1-c)^d}\sum_{k=1}^{n-1} P_{k-1}^dt_k^d.\]
Since the homothety $\mathcal{H}_{\mathbf{o}}^{P_{n-1}}$ scales the Lebesgue's measure by the factor $P_{n-1}^d$, the second term is
\[ P_{n-1}^d \Big| \bigcup_{i_n} \big( (T_{i_n} + \mathbf{w}_{i_n}') \setminus \mathcal{H}_{\mathbf{o}}^{(1-t_n)}(T)\big)\Big|,\]
which by \ref{item:induction11} is less than $P_{n-1}^d|T| t_n^d(1-c)^{-d}$. Altogether this proves (i).

In order to prove (ii), if the family $(\Sigma_i,\mathbf{x}_i)_{i \in I}$ is affinely separated, take $i\neq j\in I^n$. If $i_n=j_n$, then $i'\neq j'$ and by \ref{item:inductionn2} 
\begin{equation*}
   \Big( \operatorname{int}(\ocone{\Sigma_i^{(n)}}{ \mathbf{v})}+\mathbf{w}_{i'}^{i_n}\Big)\cap \Big( \operatorname{int}(\ocone{\Sigma_j^{(n)}}{ \mathbf{v})}+\mathbf{w}_{j'}^{j_n}\Big)=\emptyset,
\end{equation*}
or equivalently since $\mathbf{w}_i = \mathbf{w}_{i'}^{i_n} + P_{n-1} \mathbf{w}_{i_n}'$ and $\mathbf{w}_j = \mathbf{w}_{j'}^{j_n} + P_{n-1} \mathbf{w}_{i_n}'$,
\begin{equation}\label{eq:cones_disjoint_interiors}
    \Big( \operatorname{int}(\widetilde{T}_i)+\mathbf{w}_{i}\Big)\cap \Big( \operatorname{int}(\widetilde{T}_j)+\mathbf{w}_{j}\Big)=\emptyset .
\end{equation}

If $i_n\neq j_n$, since by \ref{item:inductionn2} we have that
\begin{equation*}
    \operatorname{int}(\widetilde{T}_i)+\mathbf{w}_{i} = \Big( \operatorname{int}(\ocone{\Sigma_{i'}^{i_n}}{ \mathbf{v})}+\mathbf{w}_{i'}^{i_n}\Big)+P_{n-1}\mathbf{w}_{i_n}'
   \subset \operatorname{int}\Big(\mathcal{H}^{P_{n-1}}_\mathbf{o}(\widetilde{T}_{i_n}+\mathbf{w}_{i_n}')\Big),
\end{equation*}
and by \ref{item:induction12}, the two sets $\mathcal{H}^{P_{n-1}}_\mathbf{o}(\widetilde{T}_{i_n}+\mathbf{w}_{i_n}')$ and $\mathcal{H}^{P_{n-1}}_\mathbf{o}(\widetilde{T}_{j_n}+\mathbf{w}_{j_n}')$ have disjoint interiors, we deduce that \eqref{eq:cones_disjoint_interiors} holds also in this case.

Regardless of the fact that the partition is affinely separated, the previous containment identities also imply that for any $i\in I^n$:
\begin{equation*}
   \widetilde{T}_i+\mathbf{w}_{i} \subset \mathcal{H}^{P_{n-1}}_\mathbf{o}(\widetilde{T}_{i_n}+\mathbf{w}_{i_n}')\subset \mathcal{H}^{P_{n-1}}_\mathbf{o}( \mathcal{H}^{1-t_n}_\mathbf{o}(\widetilde T)) =  \mathcal{H}^{P_n}_\mathbf{o}(\widetilde{T}),
\end{equation*}
where the inclusions are due to \ref{item:inductionn2} and \ref{item:induction12} respectively.

(iii) is also immediate: by \ref{item:inductionn3}, we have
\begin{equation*}
    \Sigma_i^{(n)}+\mathbf{w}_i=\big(\Sigma_{i'}^{i_n}+\mathbf{w}_{i'}^{i_n}\big)+P_{n-1}\mathbf{w}_{i_n}'\subset \mathcal{H}^{P_{n-1}}_\mathbf{o}(\Sigma_{i_n})+P_{n-1}\mathbf{w}_{i_n}'=\mathcal{H}^{P_{n-1}}_\mathbf{o}(\Sigma_{i_n} + \mathbf{w}_{i_n}'),
\end{equation*}
which by \ref{item:induction13} is contained in $\mathcal{H}^{P_{n-1}}_\mathbf{o}(\mathcal{H}^{1-t_n}_\mathbf{o}(\Sigma))=\mathcal{H}^{P_n}_\mathbf{o}(\Sigma)$.
\end{proof}

\end{document}
